\documentclass[../main.tex]{subfiles}

\begin{document}

\addcontentsline{toc}{section}{Preface}
\section*{Preface}

These notes started as a study guide for CS 111 at UC Santa Barbara and turned into a habit. Whenever I ran into an idea I couldn't explain from scratch, I wrote it down until I could. Feynman's test (if you can't explain something simply, you don't really understand it) turns out to be a reliable way to find gaps, and research kept turning up new ones in eigenvalue algorithms, iterative solvers, spectral graph methods, and Koopman theory. What you're reading is a working reference I wrote as a graduate student, mostly for my own use.

\medskip

Computers were first built for numerical work. A ``computer'' was originally a person, someone paid to evaluate functions by hand. During the Second World War the U.S. Army employed rooms full of them to produce artillery firing tables, and the backlog was large enough to justify building \href{https://www.computerhistory.org/revolution/birth-of-the-computer/4/78/323}{ENIAC} (1945). Among the machine's first jobs were calculations for thermonuclear weapons design, work in which John von Neumann was closely involved. His \href{https://projecteuclid.org/journals/bulletin-of-the-american-mathematical-society-new-series/volume-53/issue-11/Numerical-inverting-of-matrices-of-high-order/bams/1183511222.full}{1947 paper with Goldstine} laid the foundations of numerical stability analysis. In 1950, Jule Charney's group used ENIAC to run the \href{https://maths.ucd.ie/~plynch/eniac/CFvN-1950.pdf}{first computer weather forecast}, a calculation \href{https://archive.org/details/weatherpredictio00richrich}{Lewis Fry Richardson} had once estimated would take 64,000 people working in concert. Later, the \href{https://ntrs.nasa.gov/citations/19760016247}{Apollo Guidance Computer} navigated to the Moon with a few kilobytes of working memory, and the \href{https://www.ams.org/journals/mcom/1965-19-090/S0025-5718-1965-0178586-1/S0025-5718-1965-0178586-1.pdf}{Cooley--Tukey FFT} (1965) cut a core signal-processing computation from $O(n^2)$ to $O(n \log n)$.

\medskip

Most of what follows is linear algebra. Nearly every numerical method works by reducing a nonlinear problem to a sequence of linear ones, because linear problems are the ones we understand best. We know when solutions exist, how sensitive they are to perturbations, and how to compute them reliably. Newton's method, finite element discretization, PCA, and gradient-based optimization all come down to repeated linearization. Once linear algebra feels familiar, most of the other methods in these notes are easier to follow.

\medskip

For a more rigorous treatment, start with \href{https://math.nyu.edu/~shelley/Classes/SciComp/BindelGoodman.pdf}{Principles of Scientific Computing} by Bindel and Goodman and G. W. Stewart's \href{https://epubs.siam.org/doi/book/10.1137/1.9781611971491}{Afternotes on Numerical Analysis}. For a broader tour of the literature, the \href{https://www.cs.cornell.edu/courses/cs6241/2025sp/reading.html}{CS 6241 reading list} at Cornell is a great place to look. Much of the perspective here comes from \textbf{CS 210A (Matrix Analysis)}, taught by Professor Shivkumar Chandrasekaran in the ECE department at UC Santa Barbara. If you get the chance to take it, do.

\vspace{1em}
\begin{flushright}
\textit{Aditya Dendukuri} \\
\textit{UC Santa Barbara, 2025}
\end{flushright}

\newpage

\end{document}
